\section{Generation, conservation and realization}
\label{nuclear:introduccion}
The object of Triadic Modular Holography is a state generated by positions and operations that preserves the information of its evolution. Numerical evaluation and boundary incidence are correlated realizations of this antecedent. Contraction of compatible cylinders determines values; prolongation of their incidences preserves relations that distinguish the generating states. This priority of construction also governs Dimensional Mechanics: action, duration, extension and mass arise through readers whose domains preserve orientation, carry and the history of operations.

This article studies a precise structural question: what is conserved when a readout reduces the state, where the complement resides, and how both value and dynamics can be reconstructed from it. A complete answer requires three levels. First, emissions and their refinement are constructed. A transport identity is then proved on the realized state. Finally, we determine what a specific reader observes and what physical consequences follow from its composition with the action and scale readers. The development follows this order in each application.

\subsection{Principal thesis and statements}
Evolutionary conservation of unity admits an operator formulation. For two transports \(B,C\) between spaces equipped with forms \(b,b'\), with
\[
B^{*}b'B=C^{*}b'C=b,
\]
the collective representation of the nine positions yields
\begin{equation}
\mathcal U(B,C)=\frac19
\begin{pmatrix}8B+C&\sqrt8(C-B)\\ \sqrt8(C-B)&B+8C\end{pmatrix},
\qquad
\mathcal U(B,C)^*(b'\oplus b')\mathcal U(B,C)=b\oplus b.
\label{nuclear:eq:rectora}
\end{equation}
The first column contains the readout and memory generated by a state with no incoming memory. The second column describes the contribution of memory that enters again. The equality preserves a positive form when \(b\) is an inner product, and preserves oriented area when \(b\) is a symplectic form. These two interpretations use explicitly different structures.

This law yields four connected consequences. The first is stable reconstruction of the complete register, including its terminal component. The second recovers from dodecaphasic memory the direction determining the dimensional projector. The third extends conservation to areas, volumes and higher exterior degrees. The fourth identifies, in a specified quantum realization, the part of memory that changes the purity of a partial readout. Each consequence follows from a specified mathematical operation on the preceding one.

\begin{figure}[htbp]
\centering
\begin{tikzpicture}[every node/.style={font=\small,align=center},b/.style={draw,rounded corners,inner sep=6pt,text width=39mm},>=Stealth]
\node[b] (g) {Emission and state\\APP--TRIT--TPK};
\node[b,right=12mm of g] (u) {Complete transport\\and conserved form};
\node[b,right=12mm of u] (l) {Readout and memory\\with reconstruction};
\node[b,below=12mm of g] (i) {Incidence and\\refinement};
\node[b,below=12mm of u] (a) {Action and\\dimensional realization};
\node[b,below=12mm of l] (s) {Covariance, spectrum\\and reduced entropy};
\draw[->] (g)--(u);\draw[->] (u)--(l);\draw[->] (g)--(i);
\draw[->] (u)--(a);\draw[->] (l)--(s);\draw[->] (i)--(a);\draw[->] (a)--(s);
\end{tikzpicture}
\caption{Dependencies between realizations. The arrows represent operations proved in the body of the article; the scalar readout preserves an antecedent richer than its own value.}
\end{figure}

\subsection{Complete memory and partial observation}
For a sequence of compressions \(T_n\) with complements \(D_n\), write
\(P_0=I\), \(P_{n+1}=T_nP_n\). The local equality
\(T_n^*T_n+D_n^*D_n=I\) implies
\[
I=A_\infty+\sum_{n\ge0}P_n^*D_n^*D_nP_n,
\qquad
A_\infty=\underset{N\to\infty}{\operatorname{s-lim}}P_N^*P_N.
\]
Reconstruction uses \(A_\infty^{1/2}v\) and all the \(D_nP_nv\). Strong convergence of the terminal Gram operator suffices for conservation; it is proved by positive monotonicity, independently of convergence of the terminal vector. The register contains amplitudes and correlations, whereas their intensities provide the norm balance. This distinction determines the recoverable information.

In a specific finite register, the incidence of twelve positions in the Witt hexads has Gram matrix
\(36I+30\mathbf1\mathbf1^*\). Its normalization preserves exactly the centred part of the register. Incorporating the mean recovers its twelve coordinates. The subsequent dimensional map uses this transported information and preserves its dependence on the ordered register \(K\).

\subsection{A calculable spectral consequence}
Let \(J^2=-I\), \(J^{\mathsf T}=-J\), and let \(H\) be symplectic in a canonical realization of dimension \(2d\). Define
\[
D=\frac{\sqrt8}{9}(H-I),\qquad
D_{\mathrm{anti}}=\frac{D+JDJ}{2},\qquad
W=\frac{8I+HH^{\mathsf T}}9.
\]
The article proves
\begin{equation}
\boxed{-(JW)^2=I+[D,J][D,J]^{\mathsf T}
=I+4D_{\mathrm{anti}}D_{\mathrm{anti}}^{\mathsf T}.}
\label{nuclear:eq:espectro}
\end{equation}
For two identical pure Gaussian inputs and transport \(\mathcal U(I,H)\), the matrix \(W\) is the normalized reduced covariance. Hence
\[
\operatorname{Tr}\rho_{\mathrm{vis}}^2
=\det\bigl(I+4D_{\mathrm{anti}}D_{\mathrm{anti}}^{\mathsf T}\bigr)^{-1/4}.
\]
The formula relates a memory operator to the purity functional of the reduced state. Global purity remains equal to one. In particular, mixing depends on the antilinear component relative to the transported complex structure, rather than on an indiscriminate identification of memory with entropy.

The choice of Gaussian representation, input state and observed subsystem forms part of the statement. The genealogy specifies the transports being composed; the protocol specifies the observation whose distribution is calculated. This specification allows reproducible comparison between realizations and separates the theorem from any subsequent experimental attribution.

\subsection{Order of exposition}
The initial chapters develop the common primitives, prolongation, numerical regions, register and exceptional incidence in full. Action readers and dimensional representations follow. The second part brings together the proofs of conservation, recovery, holonomy and reversible extension. The final part composes these results with physical actions and obtains the spectrum of reduced memory. The bibliography situates the tools of representation after the construction they realize.

Finite symmetry is expressed through actions and intertwiners. Continuous variational symmetry is expressed through an action, a variation and its conserved charge. Noether's theorem applies in this second domain; transport of exceptional incidence preserves its own algebraic content. Their relationship is established through the transported observables and forms, with their differences of type kept explicit.

% BEGIN R27_FRAMING_INTRO_EN
Preservation is also examined before linear representations are
introduced. The emitter's decimal encoding has a check congruence
and a complete alphabet of \(42\) blocks, proved through its signatures
and seeds. Incidences between emission halves join the regions while
preserving the provenance of every edge. Two records with a common
histogram and different correlation establish a concrete distinction
between distribution and order. Archive laws receive that order together
with the amplitudes and terminal data stated in their hypotheses.

Normalization of the multiplicative logarithm is distinguished from its
inverse by the least-positive-generator criterion. In integer closure,
controls of carry direction, terminal endpoint, rotation, and local
perturbation keep the remaining data fixed. Boundary selection applies
saturation, dual neutrality, and orientation successively. These operations
precede recovery readings and retain their domains, so that each inverse
acts on the data actually produced by its construction.
% END R27_FRAMING_INTRO_EN

% BEGIN R28_FRAMING_INTRO_EN
The linkage of records is made explicit through the emitter's multigraph
and its phase quotient: five minimal closures join the four incidences
required for each region. The derivation of boundary occupancy and the
parity of readers under oriented reversal complete the description of
the data received by bilateral transport and retained by its coupling
readings.
% END R28_FRAMING_INTRO_EN
